\documentclass[11pt]{article}

\usepackage[margin=1in]{geometry}
\usepackage{amsmath, amssymb}
\usepackage{booktabs}
\usepackage{listings}
\usepackage{xcolor}
\usepackage{graphicx}
\usepackage{float}

\title{World and Launch Documentation}
\author{Alexander Greus, Kenneth Thompson, Dylan Zemlin, Spencer Smith}
\date{\today}

\begin{document}

\maketitle

\part*{World File}
This file specifies bunch of information about the map that we made for this project. Most importantly it specifies the model that we created to simulate the given assignment map specification. The model section lists all the walls that make it up, their positions/'poses' and their size. It also lists the collision separately from their visuals, which though its probably not as important for our model, since its relatively simple, it probably is important to separate visual and colliding geometry to optimize their representations for each. As well it also specifies the ground plane model. We initially ran into a humorous bug when we first tried running the turtle bot on our map: we had forgotten to put a ground plane and immediately the turtle bot began plummeting into infinity. For all the models in our world file there also parameters that we didn't really mess with, such as friction or visual material. As well, there are parameters for the world that seem very interesting but that we did not change, such as gravity or magnetic field. There are also parameters for physics, and rendering plugins.

\part*{Launch File}
Our launch file ended up being pretty simple since we were able to reuse or 'include' the launch files from the other turtle bot modules. It begins with setting up the config arguments, most importantly the world file matching the world file we created. It also sets up some other turtle bot configurations, such as the base model, battery, etc. The launch file also includes a GUI argument, since we were having issues initially bringing up the GUI's for Gazebo and RVIZ initially. The launch file then starts up Gazebo with our world file. Next, we launch the gmapping package for the mapping capabilities. We then spawn our turtle bot into the gazebo world, using the arguments we discussed previously.

\end{document}
